\documentclass[11pt,a4paper]{article}

\usepackage[utf8]{inputenc}
\usepackage[T1]{fontenc}
\usepackage[margin=2.5cm]{geometry}
\usepackage{amsmath,amssymb,bm}
\usepackage{booktabs}
\usepackage{array}
\usepackage{enumerate}
\usepackage{xcolor}
\usepackage{hyperref}
\usepackage{tikz}
\usepackage{verbatim}

\hypersetup{colorlinks=true, linkcolor=blue, urlcolor=blue}

\newcommand{\points}[1]{\textbf{(#1\,\%)}}
\newcommand{\subpoints}[1]{\textit{(#1\,\%)}}
\newcommand{\sol}{\noindent\textbf{Solution}}

\title{TMA4268 Statistical Learning V2026\\[0.3em]
\large Mock Exam 4 (solutions)}
\author{Compiled for Anders Bekkevard\\[0.2em]
\small Companion to \texttt{mock-exam-4.tex} (same directory).}
\date{Mock for: May 18, 2026}

\begin{document}
\maketitle

\noindent\textit{This document is a worked-solution proposal in the style of the official
Stefanie/Sara solutions to the 2024 and 2025 finals. Point values are quoted in the
heading of each solved sub-part. Partial-credit hints are inline. Refer back to
\texttt{mock-exam-4.tex} for the problem statements.}

\vspace{0.6em}
\hrule
\vspace{1em}

%=========================================================
\section*{Problem 1 \points{10} --- Fill-in-the-blank}
%=========================================================

\sol{} \emph{(1\,P per blank.)}

\begin{enumerate}[(1)]
  \item \textbf{collinearity}
  \item \textbf{ridge regression}
  \item \textbf{$k$-fold cross-validation}
  \item \textbf{nested cross-validation}
  \item \textbf{mini-batch SGD}
  \item \textbf{backpropagation}
  \item \textbf{dropout}
  \item \textbf{early stopping}
  \item \textbf{label smoothing}
  \item \textbf{gradient boosting}
\end{enumerate}

\noindent\emph{Grading: 1\,P per correct blank. No partial credit for ``close'' alternatives
(e.g.\ ``the lasso'' for (2), where the passage explicitly says ``a quadratic penalty'').}

\vspace{0.6em}
\hrule
\vspace{1em}

%=========================================================
\section*{Problem 2 \points{28} --- Multiple choice, T/F, short numeric}
%=========================================================

\subsection*{a) Bias--variance and double descent \subpoints{3}}
\sol{} \emph{(1\,P per statement.)}
\begin{enumerate}[(i)]
  \item \textbf{True}. The bias--variance decomposition is an algebraic identity that follows from
        adding/subtracting $\mathbb{E}[\hat f(x_0)]$ inside the squared expectation; it does not
        require $\hat f$ to be the squared-error minimiser.
  \item \textbf{False}. The decomposition continues to hold in the over-parameterised regime; what
        ``double descent'' shows is that the \emph{variance} term can decrease again past the
        interpolation threshold, not that the identity itself breaks.
  \item \textbf{False}. This is the canonical reason the prof calls the result a \emph{decomposition}
        rather than a \emph{trade-off}. Regularisation, more data, or structural changes can reduce
        variance \emph{without} a matching increase in bias.
\end{enumerate}

\subsection*{b) Cross-validation and the wrong-way CV trap \subpoints{4}}
\sol{} \emph{(1\,P per statement.)}
\begin{enumerate}[(i)]
  \item \textbf{True}. With $k=10$ the model is fit exactly $10$ times (once per held-out fold). LOOCV
        is the $k=n$ special case.
  \item \textbf{False}. This is the classic \emph{wrong-way CV} trap (ISLR \S5.3 / lecture 6).
        Pre-selecting predictors on the full data leaks $y$-information into every fold, so CV
        underestimates the test error. The selection step must move \emph{inside} each fold.
  \item \textbf{True}. This is exactly the structural definition of nested CV: inner loop selects,
        outer loop assesses the procedure that includes the selection step.
  \item \textbf{True}. PRESS identity: $\mathrm{LOOCV} = \frac{1}{n}\sum_i\!\left(\frac{y_i-\hat y_i}{1-h_{ii}}\right)^2$,
        so OLS-LOOCV needs only one full fit plus the diagonal of the hat matrix.
\end{enumerate}

\subsection*{c) Neural network parameters and forward pass \subpoints{3}}

\sol{} \subpoints{1} \textbf{(i)} Each layer contributes (\#inputs\,$+\,1$)\,$\times$\,(\#outputs)
parameters; the $+1$ counts the bias.
\begin{align*}
\text{input}(4)\to\text{hidden}(5): &\quad 4\cdot 5 + 5 = 25, \\
\text{hidden}(5)\to\text{hidden}(3): &\quad 5\cdot 3 + 3 = 18, \\
\text{hidden}(3)\to\text{output}(1): &\quad 3\cdot 1 + 1 = 4, \\
\text{Total:} &\quad \boxed{25 + 18 + 4 = 47} \text{ parameters.}
\end{align*}

\medskip
\sol{} \subpoints{2} \textbf{(ii)} Pre-activation:
$$z \;=\; b + \sum_{j=1}^4 w_j x_j
   \;=\; -1.5 + 1.0\cdot 1 + (-0.5)\cdot 2 + 2.0\cdot(-1) + 0\cdot 4$$
$$\phantom{z} \;=\; -1.5 + 1 - 1 - 2 + 0 \;=\; -3.5.$$
ReLU: $\max(0,\,-3.5) = \boxed{0}$.

\noindent\emph{Grading: 1\,P for $z=-3.5$, 1\,P for the ReLU step. Common slip: forgetting the
sign on the bias (gives $z=-0.5$, still ReLU $=0$); accept full credit only if the bias term
is correctly placed.}

\subsection*{d) Backprop and mini-batch SGD \subpoints{3}}
\sol{} \emph{(0.75\,P per statement.)}
\begin{enumerate}[(i)]
  \item \textbf{False}. Backpropagation is the \emph{algorithm that computes the gradient} (via the
        chain rule). The optimiser that uses that gradient to update weights is (mini-batch) SGD.
  \item \textbf{True}. A uniformly drawn mini-batch gives an unbiased estimator of the full-data
        gradient: $\mathbb{E}_{\text{batch}}[\nabla L_{\text{batch}}] = \nabla L_{\text{full}}$.
  \item \textbf{True}. Larger batches have lower-variance gradient estimates; the gradient noise
        from small batches is itself a (mild) regulariser, so reducing it reduces that regularisation.
  \item \textbf{True}. With $L = \tfrac{1}{2}(y_i-\hat f(x_i))^2$, $\partial L/\partial \hat f = -(y_i-\hat f(x_i))$;
        this is the seed of the backward pass.
\end{enumerate}

\subsection*{e) Neural-network regularization \subpoints{3}}
\sol{} \emph{(0.75\,P per statement.)}
\begin{enumerate}[(i)]
  \item \textbf{False}. Dropout is active \emph{only} at training time. At test time all units are
        used, with outputs scaled appropriately (or weights scaled, in ``inverted dropout'').
  \item \textbf{False}. $50\%$ is too aggressive for fully-connected hidden layers in this course;
        the prof's typical default is $\sim 20\%$. Rates near $0.5$ tend to under-fit on small
        tabular data.
  \item \textbf{False}. Early stopping monitors the \emph{validation} loss and returns the
        parameters from the epoch that minimised it. Stopping when training loss plateaus would
        defeat the purpose.
  \item \textbf{True}. Label smoothing replaces one-hot targets with softened ones; the prof
        explicitly motivated it by label noise in the training data.
\end{enumerate}

\subsection*{f) Boosting flavors \subpoints{4}}
\sol{} \emph{(1\,P per statement.)}
\begin{enumerate}[(i)]
  \item \textbf{True}. AdaBoost up-weights misclassified observations each round via
        $w_i \leftarrow w_i e^{\alpha_m \mathbb{1}[y_i\neq G_m(x_i)]}$.
  \item \textbf{True}. For squared-error loss $L=\tfrac{1}{2}(y-f)^2$, the negative gradient w.r.t.\
        $f$ is $-(\partial L/\partial f) = y-f =$ residual. Fitting the next tree to residuals is
        therefore gradient descent in function space.
  \item \textbf{False}. The relationship is the opposite: smaller $\nu$ means each tree contributes
        \emph{less}, so \emph{more} trees are needed to reach the same fit. (Roughly: halving $\nu$
        doubles the required $M$.)
  \item \textbf{True}. XGBoost's distinguishing features are (a) second-order Taylor expansion of
        the loss (uses both gradient and Hessian), and (b) explicit $L_1/L_2$ regularisation on the
        leaf weights, in addition to engineering optimisations.
\end{enumerate}

\subsection*{g) Collinearity \subpoints{3}}
\sol{} \emph{(1\,P per statement.)}
\begin{enumerate}[(i)]
  \item \textbf{True}. With near-singular $\mathbf X^\top\mathbf X$, the diagonal entries of
        $\sigma^2(\mathbf X^\top\mathbf X)^{-1}$ blow up, inflating standard errors. The joint
        $F$-test on the collinear bloc can still be highly significant.
  \item \textbf{False}. Collinearity inflates the variance of individual \emph{coefficient
        estimates}, but the \emph{prediction} $\hat y = \mathbf X\hat{\boldsymbol\beta}$ can still
        be accurate (the unstable directions are exactly the ones that don't affect the fitted
        values much).
  \item \textbf{False}. \emph{Perfect} collinearity ($\texttt{age\_decades} = \texttt{age}/10$) makes
        $\mathbf X^\top\mathbf X$ exactly singular: the OLS estimates are \emph{not defined} at all
        (no unique solution), so it is not the case that the estimates exist but with infinite SEs.
        Most software silently drops one of the two columns.
\end{enumerate}

\subsection*{h) Principal component analysis \subpoints{3}}

\sol{} \subpoints{1} \textbf{(i)} For standardised variables each has variance $1$, so the total
variance equals the number of variables: $\sum_j \lambda_j = p = 7$
($2.5+1.5+1.1+0.8+0.5+0.4+0.2 = 7$, \checkmark). Cumulative proportion of explained variance:
\begin{align*}
\text{PC1:}\ & 2.5/7 = 0.357, \\
\text{PC1--2:}\ & 4.0/7 = 0.571, \\
\text{PC1--3:}\ & 5.1/7 = 0.729, \\
\text{PC1--4:}\ & 5.9/7 = 0.843 \;\geq\; 0.80. \;\checkmark
\end{align*}
So $\boxed{4}$ principal components are needed to reach at least $80\%$.

\medskip
\sol{} \subpoints{1} \textbf{(ii)} Score $z^*_1 = \phi_1^\top x^*$:
\begin{align*}
z^*_1 &= 0.50\cdot 2 + (-0.40)\cdot 0 + 0.45\cdot 1 + 0.30\cdot(-1) + (-0.35)\cdot 0 + 0.40\cdot 1 + 0.15\cdot 2 \\
      &= 1.00 + 0 + 0.45 - 0.30 + 0 + 0.40 + 0.30 \;=\; \boxed{1.85}.
\end{align*}

\medskip
\sol{} \subpoints{1} \textbf{(iii)} \textbf{False}. PCA is performed on the (centered)
\emph{covariance matrix} of the predictors, which is sensitive to scale. Without standardisation,
variables measured on larger scales dominate the first PC; standardising removes this scale
dependence and is the default whenever the variables are on different scales.

\subsection*{i) Hierarchical clustering and $K$-means \subpoints{2}}
\sol{} \emph{(1\,P per statement.)}
\begin{enumerate}[(i)]
  \item \textbf{True}. At the first merge each observation is its own cluster, so all three
        linkages (single / complete / average) reduce to ``find the smallest off-diagonal entry of
        $D$.'' The linkage rule only kicks in from the second merge onward (when at least one
        cluster contains more than one point).
  \item \textbf{True}. The standard $K$-means algorithm (Lloyd's algorithm) is only guaranteed to
        converge to a \emph{local} minimum of the within-cluster sum of squares; different random
        initialisations can yield different solutions. The recommended practical fix is exactly
        what the statement says: run the algorithm from several random starts and keep the one
        with the smallest within-cluster SS.
\end{enumerate}

\vspace{0.6em}
\hrule
\vspace{1em}

%=========================================================
\section*{Problem 3 \points{16} --- Theory and hand calculations}
%=========================================================

\subsection*{a) The mathy one --- the bias--variance decomposition \subpoints{8}}

\sol{} \subpoints{2} \textbf{(i)} Substitute $y_0 = f(x_0) + \varepsilon$ and expand:
\begin{align*}
\mathbb{E}\!\left[(y_0-\hat f(x_0))^2\right]
  &= \mathbb{E}\!\left[(f(x_0) + \varepsilon - \hat f(x_0))^2\right] \\
  &= \mathbb{E}\!\left[(f(x_0)-\hat f(x_0))^2\right]
     + 2\,\mathbb{E}\!\left[\varepsilon\bigl(f(x_0)-\hat f(x_0)\bigr)\right]
     + \mathbb{E}[\varepsilon^2].
\end{align*}
The cross term \emph{vanishes}: $\varepsilon$ is independent of the training set $\mathcal{D}$
(hence of $\hat f$) and of the fixed $f(x_0)$, and $\mathbb{E}[\varepsilon]=0$, so
$$\mathbb{E}\!\left[\varepsilon\bigl(f(x_0)-\hat f(x_0)\bigr)\right]
   = \mathbb{E}[\varepsilon]\cdot \mathbb{E}\!\left[f(x_0)-\hat f(x_0)\right] = 0.$$
And $\mathbb{E}[\varepsilon^2] = \mathrm{Var}(\varepsilon) = \sigma^2$ since $\mathbb{E}[\varepsilon]=0$.
Therefore
$$\boxed{\;\mathbb{E}[(y_0-\hat f(x_0))^2] = \mathbb{E}[(f(x_0)-\hat f(x_0))^2] + \sigma^2.\;}$$

\noindent\emph{Grading: 1\,P for the expansion; 1\,P for explicitly killing the cross term using
independence + zero mean. Deduct 0.5 if the cross-term is dropped without justification.}

\medskip
\sol{} \subpoints{3} \textbf{(ii)} Add and subtract $\mathbb{E}[\hat f(x_0)]$ inside the squared
reducible term:
\begin{align*}
\mathbb{E}\!\left[(f(x_0)-\hat f(x_0))^2\right]
&= \mathbb{E}\!\left[\bigl((f(x_0)-\mathbb{E}[\hat f(x_0)]) + (\mathbb{E}[\hat f(x_0)] - \hat f(x_0))\bigr)^2\right] \\
&= \bigl(f(x_0)-\mathbb{E}[\hat f(x_0)]\bigr)^2 \\
&\quad + 2\bigl(f(x_0)-\mathbb{E}[\hat f(x_0)]\bigr)\,\mathbb{E}\!\left[\mathbb{E}[\hat f(x_0)] - \hat f(x_0)\right] \\
&\quad + \mathbb{E}\!\left[(\hat f(x_0) - \mathbb{E}[\hat f(x_0)])^2\right].
\end{align*}
The first term is constant (pulled out of the expectation since $f(x_0)$ is fixed and
$\mathbb{E}[\hat f(x_0)]$ is a constant). The middle (cross) term vanishes:
$$\mathbb{E}\!\left[\mathbb{E}[\hat f(x_0)] - \hat f(x_0)\right] = \mathbb{E}[\hat f(x_0)] - \mathbb{E}[\hat f(x_0)] = 0.$$
The third term is exactly $\mathrm{Var}(\hat f(x_0))$ by definition. Hence
$$\boxed{\;\mathbb{E}\!\left[(f(x_0)-\hat f(x_0))^2\right]
   = \underbrace{\bigl(f(x_0)-\mathbb{E}[\hat f(x_0)]\bigr)^2}_{\mathrm{Bias}^2[\hat f(x_0)]}
   + \underbrace{\mathrm{Var}(\hat f(x_0))}_{\text{variance}}.\;}$$

\noindent\emph{Grading: 1\,P for the add-subtract trick; 1\,P for vanishing of the cross term; 1\,P
for identifying the two surviving terms as $\mathrm{Bias}^2$ and $\mathrm{Var}$.}

\medskip
\sol{} \subpoints{1} \textbf{(iii)} Combining (i) and (ii),
$$\mathbb{E}\!\left[(y_0-\hat f(x_0))^2\right]
   = \underbrace{\bigl(f(x_0)-\mathbb{E}[\hat f(x_0)]\bigr)^2}_{\mathrm{Bias}^2}
   + \underbrace{\mathrm{Var}(\hat f(x_0))}_{\text{variance}}
   + \underbrace{\sigma^2}_{\text{irreducible}}.$$
$\sigma^2$ is the \emph{irreducible} error (the variance of the noise $\varepsilon$). The inner
$\mathbb{E}[\hat f(x_0)]$ and $\mathrm{Var}[\hat f(x_0)]$ are taken over the random training set
$\mathcal{D}$; the outer $\mathbb{E}[\cdot]$ above is jointly over $\mathcal{D}$ and over the test-point
noise $\varepsilon$.

\medskip
\sol{} \subpoints{1} \textbf{(iv)} ``Decomposition'' is preferred because the identity is an exact
algebraic equality, not a constraint that forces bias to grow when variance shrinks. Regularisers
such as ridge / lasso / dropout can reduce variance \emph{without} a matching increase in bias,
and in the over-parameterised / double-descent regime variance can even decrease further past the
interpolation threshold.

\medskip
\sol{} \subpoints{1} \textbf{(v)} At a small positive $\lambda$, lasso introduces a little
\emph{bias} (shrinkage toward zero) but reduces the \emph{variance} of the fit; the irreducible
$\sigma^2$ is unchanged. The net is that variance falls more than bias$^2$ rises, so test MSE
drops --- exactly the bias--variance lever the penalty is designed to pull.

\subsection*{b) Pseudocode --- nested $k$-fold cross-validation \subpoints{4}}

\sol{} \subpoints{3} \textbf{(i)}

{\small
\begin{verbatim}
INPUT: dataset (X, y) of size n;
       hyperparameter grid Lambda = {lambda_1, ..., lambda_L};
       K_outer = 5, K_inner = 5.

partition (X, y) into K_outer outer folds F_1, ..., F_{K_outer}

FOR k = 1 to K_outer:                                    # OUTER (assessment) loop
    test_k       = F_k                                   # held out for HONEST assessment
    trainval_k   = union of all other outer folds

    partition trainval_k into K_inner inner folds G_1, ..., G_{K_inner}

    FOR each lambda in Lambda:                           # HYPERPARAMETER selection
        FOR j = 1 to K_inner:                            # INNER (selection) loop
            val_j   = G_j
            train_j = trainval_k minus G_j
            fit model_{lambda, j} on train_j             # data used to FIT
            err_{lambda, j} = error(model_{lambda, j}, val_j)
        end
        innerCV(lambda) = mean over j of err_{lambda, j}
    end

    lambda_k* = argmin_lambda innerCV(lambda)            # selection step
    refit model_k on ALL of trainval_k at lambda_k*       # refit on full outer-train fold
    score_k    = error(model_k, test_k)                  # OUTER assessment on test_k

end

test_error_estimate = mean over k of score_k             # aggregate outer scores
RETURN test_error_estimate
\end{verbatim}
}

\noindent The structure exposes the three data uses cleanly:
\emph{fit} = \texttt{train\_j}; \emph{select} = \texttt{val\_j} (through the inner-CV
average); \emph{assess} = \texttt{test\_k} (never touched until selection is complete on
\texttt{trainval\_k}).

\noindent\emph{Grading: 1\,P for the correct nesting (outer loop wraps the inner CV); 1\,P for
labelling the three data roles (fit / select / assess); 1\,P for aggregating outer-fold scores
into the final estimate. Deduct 1 if $\lambda$ is selected once globally rather than once per
outer fold.}

\medskip
\sol{} \subpoints{1} \textbf{(ii)} The shortcut reuses the same data both for \emph{choosing}
$\lambda$ and for \emph{reporting} the test error. The argmin of CV error is itself a noisy,
optimistic statistic, so the reported error is biased \emph{downward} (a ``winner's curse'' /
selection bias) --- the very phenomenon nested CV is designed to correct.

\subsection*{c) Bootstrap standard error \subpoints{4}}

\sol{} \subpoints{2} \textbf{(i)}

{\small
\begin{verbatim}
INPUT: i.i.d. sample x = (x_1, ..., x_n); statistic theta_hat = median(x);
       number of bootstrap replicates B (e.g. B = 1000).

FOR b = 1 to B:
    draw x*_b = (x*_{b,1}, ..., x*_{b,n}) by sampling n times WITH REPLACEMENT
                                                      from (x_1, ..., x_n)
    compute   theta*_b = median(x*_b)
end

theta_bar = (1/B) * sum_{b=1}^B theta*_b

SE_boot   = sqrt( (1/(B-1)) * sum_{b=1}^B (theta*_b - theta_bar)^2 )

RETURN SE_boot
\end{verbatim}
}

\noindent The three required ingredients are explicit: (a) each bootstrap sample is drawn
\emph{with replacement} from the original sample at size $n$; (b) the recomputed statistic on each
bootstrap sample is the sample median; (c) the final SE is the sample standard deviation of the
$B$ bootstrap medians (divisor $B-1$).

\medskip
\sol{} \subpoints{1} \textbf{(ii)} For one bootstrap sample of size $n$, the probability that
observation $X_1$ is \emph{not} chosen on any given draw is $(1-1/n)^n$, so
$$P(X_1 \in \text{bootstrap sample}) = 1 - \left(1 - \tfrac{1}{n}\right)^n.$$
At $n=6$:
$$P = 1 - (5/6)^6 = 1 - 0.33490 = \boxed{0.665}\ \text{(three decimals)}.$$
As $n\to\infty$, $(1-1/n)^n \to 1/e$, so
$$P(X_1 \in \text{boot}) \;\longrightarrow\; 1 - 1/e \;\approx\; \boxed{0.632}.$$

\medskip
\sol{} \subpoints{1} \textbf{(iii)} \textbf{Percentile bootstrap CI.} Compute the $B$ bootstrap
replicates $\hat\theta^*_1,\dots,\hat\theta^*_B$ as in (i); the $95\%$ percentile CI is the
interval $[\hat\theta^*_{(2.5\%)},\,\hat\theta^*_{(97.5\%)}]$, i.e.\ the $2.5$th and $97.5$th
empirical quantiles of the bootstrap distribution.

\vspace{0.6em}
\hrule
\vspace{1em}

%=========================================================
\section*{Problem 4 \points{22} --- Apartment-rent regression}
%=========================================================

\subsection*{a) OLS with collinearity, interaction, and a polynomial term \subpoints{10}}

\sol{} \subpoints{1} \textbf{(i)} Count the rows of the coefficient table: intercept,
\texttt{area}, \texttt{rooms}, \texttt{area\_per\_room}, \texttt{age}, $I(\texttt{age}^2)$,
\texttt{floor}, \texttt{distance}, \texttt{has\_balcony}, \texttt{type\_studio},
\texttt{type\_loft}, \texttt{distance:has\_balcony} $=\boxed{12}$ parameters. Consistency check:
residual d.f.\ $=n_\text{train}-p = 500-12 = 488$, matching the printout. \checkmark

\medskip
\sol{} \subpoints{2} \textbf{(ii)} The triple $(\texttt{area},\,\texttt{rooms},\,\texttt{area\_per\_room})$
satisfies $\texttt{area} = \texttt{rooms}\cdot\texttt{area\_per\_room}$, and even after
standardisation the three columns are highly correlated. This is \textbf{(near-)collinearity}:
$\mathbf X^\top\mathbf X$ is nearly singular, so the diagonal of $\sigma^2(\mathbf X^\top\mathbf X)^{-1}$
blows up and the individual standard errors of \texttt{rooms} ($1.95$) and
\texttt{area\_per\_room} ($1.40$) are large enough to wipe out their $t$-statistics, even though
the joint information they carry about size is real (\texttt{area} survives with $t=2.76$).
\textbf{Effect of dropping} \texttt{area\_per\_room}: the remaining two predictors are now less
collinear, so their SEs would \emph{shrink} and their $t$-values would (typically) become
significant.

\medskip
\sol{} \subpoints{1} \textbf{(iii)} The high $p$-values on \texttt{rooms} and
\texttt{area\_per\_room} are an artefact of inflated standard errors caused by their collinearity
with \texttt{area}, not evidence that the two predictors carry no information. A joint $F$-test on
the three size-related predictors would (very likely) reject; their information \emph{is} captured
by the model, just smeared across three highly correlated columns.

\medskip
\sol{} \subpoints{2} \textbf{(iv)} The partial effect of \texttt{age} on \texttt{rent} (holding
all other predictors fixed) is
$$\frac{\partial\,\widehat{\texttt{rent}}}{\partial\,\texttt{age}}
   \;=\; \hat\beta_{\texttt{age}} + 2\,\hat\beta_{\texttt{age}^2}\cdot \texttt{age}
   \;=\; -0.080 + 2\cdot 0.00045\cdot\texttt{age}
   \;=\; -0.080 + 0.0009\cdot\texttt{age}.$$
Setting the derivative to zero (and noting $\hat\beta_{\texttt{age}^2} > 0$ makes this a minimum):
$$\texttt{age}^* \;=\; \frac{0.080}{0.0009} \;\approx\; \boxed{88.9\ \text{years}}.$$
\emph{Interpretation:} the model's predicted rent decreases with age up to $\sim 89$ years, then
turns upward (presumably picking up classic / heritage stock).

\medskip
\sol{} \subpoints{2} \textbf{(v)} \textbf{Convention.} \texttt{distance} is standardised (mean 0,
SD 1) as stated in the preamble, so ``a 1-SD increase'' means a change of $\Delta x = 1$ in the
standardised \texttt{distance}.
\begin{align*}
\text{No balcony } (\texttt{has\_balcony}=0): &\quad
   \Delta\widehat{\texttt{rent}} = \hat\beta_{\texttt{dist}}\cdot 1 = \boxed{-1.95}\ (\text{1000 NOK}). \\
\text{With balcony } (\texttt{has\_balcony}=1): &\quad
   \Delta\widehat{\texttt{rent}} = (\hat\beta_{\texttt{dist}} + \hat\beta_{\text{dist}:\text{bal}})\cdot 1
      = -1.95 + 0.55 = \boxed{-1.40}\ (\text{1000 NOK}).
\end{align*}
Both effects are negative (apartments further from the centre rent for less), but the balcony
\emph{attenuates} the penalty by $+0.55$ --- plausibly because balconies are more valuable in the
quieter, leafier zones away from the centre, so the rent loss with distance is dampened for
balcony-equipped flats.

\medskip
\sol{} \subpoints{1} \textbf{(vi)} \textbf{Training $R^2$:} can only \emph{increase} (or stay equal)
--- adding any new column to an OLS fit cannot raise the RSS, since the smaller model is nested
inside the larger one. \textbf{Adjusted $R^2$:} depends on the size of the increase relative to the
$d.f.$ cost; if the cubic term explains real curvature it rises, otherwise it may fall.

\medskip
\sol{} \subpoints{1} \textbf{(vii)} Use an \textbf{$F$-test} (partial / nested) comparing the full
model to the reduced model that drops both \texttt{type} dummies. The null hypothesis is
$$H_0: \beta_{\texttt{type\_studio}} = \beta_{\texttt{type\_loft}} = 0,$$
i.e.\ apartment \emph{type} carries no joint information about rent after controlling for the
other predictors.

\subsection*{b) Lasso with $10$-fold CV \subpoints{5}}

\sol{} \subpoints{1} \textbf{(i)} The lasso objective on the standardised design
$\mathbf X\in\mathbb R^{n\times p}$ is
$$\hat{\boldsymbol\beta}^{\text{lasso}}
   \;=\; \arg\min_{\beta_0,\boldsymbol\beta}\;
   \frac{1}{2n}\sum_{i=1}^n\bigl(y_i - \beta_0 - \mathbf x_i^\top\boldsymbol\beta\bigr)^2
   \;+\; \lambda\sum_{j=1}^p |\beta_j|.$$
The intercept $\beta_0$ is \emph{not} penalised (otherwise the solution depends on the units of $y$);
the penalty sums absolute values of the slope coefficients only.

\medskip
\sol{} \subpoints{1} \textbf{(ii)} The $\ell_1$ ball $\{\boldsymbol\beta : \|\boldsymbol\beta\|_1 \leq t\}$
has \emph{corners on the coordinate axes}; the RSS contour first hits this ball at a corner with
positive probability, producing exact zeros. The ridge $\ell_2$ ball is smooth and round, so the
contour generically hits it at an interior point and no coefficient is set exactly to zero.

\medskip
\sol{} \subpoints{2} \textbf{(iii)} Lasso-$\hat\lambda_{\min}$ improves on OLS by $3.85 \to 3.40$
(test MSE) while zeroing $3$ coefficients. In bias--variance terms, OLS is suffering from
\emph{high variance} on the collinear $(\texttt{area},\,\texttt{rooms},\,\texttt{area\_per\_room})$
bloc identified in (a)(ii); the lasso buys back variance by introducing some bias (shrinkage and
selection), and the variance saving dominates. This is exactly the bias--variance lever the lasso
is designed to pull, and the gain confirms that part of OLS's error was variance from the
collinear directions.

\medskip
\sol{} \subpoints{1} \textbf{(iv)} \textbf{One-standard-error (1-SE) rule.} Among all $\lambda$
whose CV-MSE is within one standard error of the minimum CV-MSE, choose the \emph{largest}
(simplest / most-regularised model). The practitioner's preference is reasonable because the CV
minimum is itself a noisy estimate, $\hat\lambda_{1\text{SE}}$ yields a sparser, more
interpretable, more stable model whose test performance ($3.50$) is statistically
indistinguishable from $\hat\lambda_{\min}$ ($3.40$).

\subsection*{c) Gradient boosting \subpoints{4}}

\sol{} \subpoints{2} \textbf{(i)} One round of gradient boosting for squared-error regression,
starting from a running ensemble $\hat f^{(m-1)}$:

{\small
\begin{verbatim}
INPUT: training data (x_i, y_i), i = 1..n;
       current ensemble F_{m-1}(x);
       shrinkage nu in (0, 1]; tree size d.

# (i) compute the (negative) gradient = residual under squared loss
FOR i = 1..n:
    r_i = y_i - F_{m-1}(x_i)        # = -dL/df  with L = (1/2)(y - f)^2
end

# (ii) fit a regression tree of size d to the residuals
fit tree T_m to (x_i, r_i), i = 1..n,  with at most d splits / depth d

# (iii) update the ensemble with shrinkage nu
F_m(x)  <-  F_{m-1}(x) + nu * T_m(x)

RETURN F_m
\end{verbatim}
}

\noindent The new tree $T_m$ is fit to the \emph{current residuals}, $\nu$ scales each tree's
contribution (smaller $\nu$ = slower learning, needs larger $M$), and the running ensemble
accumulates the shrunken trees additively.

\medskip
\sol{} \subpoints{1} \textbf{(ii)} For $L(y,f) = \tfrac{1}{2}(y-f)^2$,
$$\frac{\partial L}{\partial f} = -(y-f), \qquad
   -\frac{\partial L}{\partial f} = y - f = \text{residual}.$$
So fitting the next tree to the residuals \emph{is} fitting it to the negative gradient of the
squared-error loss --- gradient descent in function space. \checkmark

\medskip
\sol{} \subpoints{1} \textbf{(iii)} Boosting's strong improvement (MSE $2.20$ vs.\ lasso $3.40$ vs.\
OLS $3.85$) suggests the rent surface is \emph{nonlinear and/or has interactions} that the linear
models cannot capture (e.g.\ nonlinear age effects, distance$\times$balcony curvature, type-specific
slopes). For interpretability, produce a \textbf{variable-importance plot} (mean reduction in
node-split loss across all $M=800$ trees), and optionally a \textbf{partial dependence plot} for
the top variables to recover direction-of-effect.

\subsection*{d) GAM degrees of freedom \subpoints{3}}

\sol{} \subpoints{1} \textbf{(i)} A cubic regression spline with $K$ interior knots uses $K+4$
basis functions including the intercept, or $K+3$ once the global intercept is removed. Here
$K=3$ (knots at $2, 4, 7$), so
$$\mathrm{bs}(\texttt{distance})\ \text{consumes}\ \boxed{K+3 = 3+3 = 6}\ \text{d.f.}$$

\medskip
\sol{} \subpoints{1} \textbf{(ii)} Total d.f.\ including the global intercept:
\begin{align*}
\text{intercept}\ &+ s(\texttt{area}, 4) + s(\texttt{age}, 5) + \mathrm{bs}(\texttt{distance}) \\
&+ \texttt{floor} + \texttt{has\_balcony} + g(\texttt{type}) \\
&= 1 + 4 + 5 + 6 + 1 + 1 + 2 = \boxed{20}\ \text{d.f.}
\end{align*}
($g(\texttt{type})$ contributes $2$ because the 3-level categorical is coded by $2$ dummy
contrasts against the \emph{flat} reference.)

\medskip
\sol{} \subpoints{1} \textbf{(iii)} The GAM is \emph{additive} by construction: each predictor
enters through its own one-dimensional smooth, so the model cannot represent \emph{interactions}
between predictors unless they are added by hand (e.g.\ \texttt{distance}:\texttt{has\_balcony},
or tensor-product smooths). Gradient-boosted trees of depth $d>1$ pick up such interactions
automatically.

\vspace{0.6em}
\hrule
\vspace{1em}

%=========================================================
\section*{Problem 5 \points{24} --- Skin-lesion classification}
%=========================================================

\subsection*{a) Logistic regression with an interaction \subpoints{9}}

\sol{} \subpoints{2} \textbf{(i)} With the interaction $\texttt{size\_mm}\!:\!\texttt{age}$, the
partial log-odds change per extra mm of \texttt{size\_mm}, at fixed \texttt{age}, is
$$\frac{\partial \hat\eta}{\partial\,\texttt{size\_mm}}
   \;=\; \hat\beta_{\texttt{size\_mm}} + \hat\beta_{\texttt{size}:\texttt{age}}\cdot \texttt{age}
   \;=\; 0.20 + 0.0020\cdot\texttt{age}.$$
The corresponding odds factor is $\exp(\cdot)$:
\begin{align*}
\texttt{age} = 30: &\quad \exp(0.20 + 0.0020\cdot 30) = \exp(0.26) \approx \boxed{1.297}, \\
\texttt{age} = 70: &\quad \exp(0.20 + 0.0020\cdot 70) = \exp(0.34) \approx \boxed{1.405}.
\end{align*}

\noindent\emph{Grading: 1\,P per age. Deduct 1\,P if both factors are $\exp(0.20)\approx 1.22$,
which ignores the interaction --- the standard trap.}

\medskip
\sol{} \subpoints{1} \textbf{(ii)} The two factors differ because the interaction term makes
the size effect itself depend on \texttt{age}: each extra year of \texttt{age} adds $0.0020$ to
the size slope. The main-effect coefficient $\hat\beta_{\texttt{size\_mm}} = 0.20$ is \emph{only}
the size effect at $\texttt{age}=0$ (an extrapolation outside the data), so it cannot summarise
the size effect on its own.

\medskip
\sol{} \subpoints{3} \textbf{(iii)} Compute the linear predictor term by term:
\begin{align*}
\hat\eta &= -6.20
   + 0.20\cdot 8
   + 2.50\cdot 0.6
   + (-0.012)\cdot 120
   + 0.030\cdot 60
   + 0.18\cdot 1
   + 0.55\cdot 1
   + 0.0020\cdot (8\cdot 60) \\
&= -6.20 + 1.60 + 1.50 - 1.44 + 1.80 + 0.18 + 0.55 + 0.96 \\
&= -6.20 + 5.15 \;=\; \boxed{-1.05}.
\end{align*}
Sigmoid step:
$$\hat p \;=\; \frac{1}{1 + e^{-\hat\eta}} \;=\; \frac{1}{1 + e^{1.05}}
   \;=\; \frac{1}{1 + 2.858} \;\approx\; \boxed{0.259}.$$

\noindent\emph{Grading: 1\,P for the per-term sum, 1\,P for $\hat\eta = -1.05$, 1\,P for the
sigmoid yielding $\hat p \approx 0.26$.}

\medskip
\sol{} \subpoints{1} \textbf{(iv)} $\hat p \approx 0.26 < 0.5$, so at the default threshold the
patient is \textbf{not flagged as malignant} (predicted class = benign).

\medskip
\sol{} \subpoints{2} \textbf{(v)}
\begin{itemize}
  \item \textbf{Sensitivity} = $P(\hat Y = \text{malignant} \mid Y = \text{malignant})$: the
        fraction of \emph{truly malignant lesions} that the model correctly flags as malignant.
        (``How many real cancers do we catch?'')
  \item \textbf{Specificity} = $P(\hat Y = \text{benign} \mid Y = \text{benign})$: the
        fraction of \emph{truly benign lesions} the model correctly leaves alone. (``How many
        healthy moles do we avoid alarming?'')
\end{itemize}
\textbf{Which to weigh more.} For a screening tool in dermatology, \emph{sensitivity} is more
important: a missed malignant melanoma is potentially fatal, whereas a false positive only
generates a follow-up biopsy. The threshold should therefore be tuned down from $0.5$ to push
sensitivity up at some cost in specificity.

\subsection*{b) Neural network with regularization \subpoints{7}}

\sol{} \subpoints{1} \textbf{(i)} Layer-by-layer:
\begin{align*}
\text{input}(6) \to \text{hidden}(32): &\quad 6\cdot 32 + 32 = 224, \\
\text{hidden}(32) \to \text{hidden}(16): &\quad 32\cdot 16 + 16 = 528, \\
\text{hidden}(16) \to \text{output}(1): &\quad 16\cdot 1 + 1 = 17, \\
\text{Total:} &\quad 224 + 528 + 17 = \boxed{769}\ \text{parameters.}
\end{align*}

\medskip
\sol{} \subpoints{2} \textbf{(ii)} \textbf{(a)} With $\sim 769$ parameters and only $4{,}500$
training images, the network is in the over-parameterised regime: $100\%$ training accuracy
combined with $0.78$ test accuracy is textbook overfitting --- the model memorises the training
set, so test accuracy degrades. Without explicit regularisation a network with this much capacity
is essentially guaranteed to overfit on a tabular signal this small.
\textbf{(b)} Mini-batch SGD with small batches injects \emph{gradient noise} into every step;
this stochastic noise acts as an implicit regulariser, biasing the optimiser toward flatter, more
generalisable minima even before dropout / early stopping / label smoothing are added.

\medskip
\sol{} \subpoints{2} \textbf{(iii)}
\begin{itemize}
  \item \textbf{Dropout ($p=0.2$).} At each training step, each hidden unit's activation is set
        to $0$ independently with probability $0.2$ (and the remaining units are typically scaled
        by $1/(1-p)$ in ``inverted dropout''). At \emph{test time} dropout is switched off and the
        full network is used --- this is the train/test asymmetry.
  \item \textbf{Early stopping.} Monitor the validation loss across epochs; halt training when the
        validation loss stops improving for a configurable patience window, and return the
        parameters from the best epoch.
  \item \textbf{Label smoothing ($\varepsilon=0.05$).} Replace the one-hot target with a softened
        version (e.g.\ $1-\varepsilon = 0.95$ on the true class, $\varepsilon = 0.05$ on the other),
        so the cross-entropy loss penalises over-confident predictions and is robust to mislabeled
        examples.
\end{itemize}

\medskip
\sol{} \subpoints{1} \textbf{(iv)} The prof's typical default is $\sim 20\%$ dropout for tabular,
relatively small networks; jumping to $50\%$ tends to \emph{under-fit} on signal-poor tabular data
(too much information is destroyed at every step) and is mostly justified on much larger CNNs with
ample training data, which is not this setting.

\medskip
\sol{} \subpoints{1} \textbf{(v)} The three regularisers traded a small increase in \emph{bias} for
a substantial decrease in \emph{variance}, lowering test error from $0.22$ to $0.14$
(accuracy $0.78\to 0.86$). The irreducible $\sigma^2$ (label noise / intrinsic difficulty) is
unchanged; the gain is pure variance reduction.

\subsection*{c) AdaBoost \subpoints{5}}

\sol{} \subpoints{2} \textbf{(i)} The exponential weight update $w_i \leftarrow w_i\exp(\alpha_m
\mathbb{1}[y_i\neq G_m(x_i)])$ inflates the weights of \emph{repeatedly misclassified}
observations: each successive weak learner is forced to focus on the hardest residual examples.
This is fundamentally different from bagging / random forests, which \emph{resample rows i.i.d.}
with uniform probability: bagging treats every observation symmetrically and averages independent
fits, while AdaBoost \emph{adapts} sequentially to the hard cases by re-weighting (no row
resampling).

\medskip
\sol{} \subpoints{1} \textbf{(ii)} If $\mathrm{err}_{m^*}=0.5$ then
$$\alpha_{m^*} = \log\!\left(\frac{1 - 0.5}{0.5}\right) = \log(1) = \boxed{0}.$$
\textbf{Effect:} the weak learner contributes \emph{nothing} to the final ensemble's prediction
($\alpha_{m^*}G_{m^*}(x) = 0$); an error rate of $0.5$ is the threshold at which a binary
classifier is no better than a coin flip and AdaBoost ignores it.

\medskip
\sol{} \subpoints{1} \textbf{(iii)} \textbf{Deploy AdaBoost.} It dominates logistic regression on
\emph{sensitivity} ($0.84$ vs.\ $0.55$) while only giving up a small amount on specificity
($0.91$ vs.\ $0.95$). For a screening tool, catching melanomas matters more than the few extra
false positives; AdaBoost flags $84\%$ of real cancers vs.\ logistic's $55\%$, a clinically
decisive gap.

\medskip
\sol{} \subpoints{1} \textbf{(iv)} \textbf{Exponential loss:} $\boxed{L(y, f) = \exp(-y\cdot f)}$
with $y\in\{-1, +1\}$. AdaBoost is gradient boosting under this loss in function space.

\subsection*{d) Class imbalance and threshold tuning \subpoints{3}}

\sol{} \subpoints{1} \textbf{(i)} If the classifier always predicts ``benign,'' it is correct on
every truly benign image and wrong on every malignant one. With $\sim 80\%$ benign in the
population, accuracy $\approx \boxed{0.80}$.

\medskip
\sol{} \subpoints{1} \textbf{(ii)} The naive classifier has sensitivity $0$ (it never flags
malignant) and specificity $1$ (it never raises a false alarm) --- useless. The logistic regression
at threshold $0.5$ achieves modest sensitivity ($\sim 0.55$ in part (c)) and high specificity
($\sim 0.95$), so even though its overall accuracy is only marginally better, it is genuinely
informative: it actually detects a majority of the cancers.

\medskip
\sol{} \subpoints{1} \textbf{(iii)} Lowering the threshold from $0.5$ to $0.3$ flags more
observations as malignant: \textbf{sensitivity increases}, \textbf{specificity decreases}, and on
the ROC curve the operating point moves \textbf{up and to the right} (higher TPR, higher FPR).

\vspace{1em}
\hrule
\vspace{0.5em}

\noindent\textbf{End of solution proposal.}\quad Total awarded: $10+28+16+22+24 = 100$ points.

\end{document}
